\documentclass[12pt]{article}
\usepackage[T1]{fontenc}
\usepackage[utf8]{inputenc}
\usepackage{lmodern}
\usepackage{textcomp}
\usepackage{enumitem}
\usepackage{geometry}
\geometry{margin=1in}
\begin{document}
\begin{center}
Answer Key - Chemistry - Undergraduate Level Assessment 17 \
\small (Answers are ordered exactly as in the question paper.)
\end{center}

\section*{Multiple Choice}
\begin{enumerate}[label=\arabic*.]
\item \textbf{Answer:} D
\\ \textit{Options:} A) Silicon; B) Diamond; C) Silicon carbide; D) Arsenic-doped silicon
\\ \textit{Note:} Arsenic-doped silicon is an n-type semiconductor because the arsenic atoms introduce extra free electrons into the silicon lattice, enhancing its conductivity by providing negative charge carriers.
\item \textbf{Answer:} D
\\ \textit{Options:} A) 6.345 x 10^-4 at 0.335 T; 9.871 x 10^-5 at 10.5 T; B) 0.793 x 10^-4 at 0.335 T; 6.931 x 10^-7 at 10.5 T; C) 1.148 x 10^-6 at 0.335 T; 3.598 x 10^-5 at 10.5 T; D) 4.126 x 10^-3 at 0.335 T; 2.142 x 10^-6 at 10.5 T
\\ \textit{Note:} The correct answer is derived from the formula for polarization in a magnetic field, which relates the magnetic moment of the proton to the strength of the magnetic field and temperature, resulting in specific calculated values of 4.126 x 10^-3 for 0.335 T and 2.142 x 10^-6 for 10.5 T at 298 K.
\item \textbf{Answer:} D
\\ \textit{Options:} A) Melting point; B) Boiling point; C) Polarizability; D) First ionization energy
\\ \textit{Note:} The first ionization energy of argon is lower than that of neon because argon has a larger atomic radius and more electron shielding, which reduces the effective nuclear charge experienced by its outermost electrons, making them easier to remove.
\item \textbf{Answer:} D
\\ \textit{Options:} A) F; B) Si; C) O; D) Ca
\\ \textit{Note:} Calcium (Ca) has the lowest electron affinity among the given atoms because it is an alkaline earth metal with a stable electron configuration that resists gaining additional electrons, resulting in a lower tendency to attract electrons compared to nonmetals.
\item \textbf{Answer:} A
\\ \textit{Options:} A) HCl; B) AgCl; C) CaCl 2; D) CCl 4
\\ \textit{Note:} HCl has the lowest melting point because it is a molecular compound with weak van der Waals forces between its molecules, compared to the stronger ionic or covalent bonds present in the other compounds.
\end{enumerate}

\section*{True / False}
\begin{enumerate}[label=\arabic*.]
\setcounter{enumi}{5}
\item \textbf{Answer:} True
\\ \textit{Options:} A) True; B) False
\\ \textit{Note:} The balanced redox equation for the reaction between permanganate ions (MnO 4-) and iodide ions (I-) shows that six moles of MnO 4- react with five moles of I-, establishing a stoichiometric ratio of 6:5.
\item \textbf{Answer:} False
\\ \textit{Options:} A) True; B) False
\\ \textit{Note:} Nuclear magnetic resonance spectroscopy is a technique that utilizes magnetic fields and radiofrequency radiation to analyze molecular structures through the behavior of atomic nuclei, rather than a light-scattering method.
\item \textbf{Answer:} False
\\ \textit{Options:} A) True; B) False
\\ \textit{Note:} The 1 H resonance of CH 3 F is shifted downfield compared to CH 3 Cl due to the greater electronegativity of fluorine, which causes a larger deshielding effect on the hydrogen nuclei.
\item \textbf{Answer:} True
\\ \textit{Options:} A) True; B) False
\\ \textit{Note:} Protons in a molecule with a rotational correlation time of 1 ns experience efficient spin-lattice relaxation at a magnetic field of 3.74 T due to optimal energy transfer mechanisms that enhance relaxation rates at this specific correlation time.
\item \textbf{Answer:} False
\\ \textit{Options:} A) True; B) False
\\ \textit{Note:} The statement is false because the number of allowed energy levels in a nucleus is determined by its angular momentum and parity, and for 55 Mn, the available nuclear states are not limited to five distinct energy levels.
\end{enumerate}

\section*{Long Form Response}
\begin{enumerate}[label=\arabic*.]
\setcounter{enumi}{10}
\item \textbf{Answer:} The ratio of line intensities in the EPR spectrum of the t-Bu radical (CH 3)3 C•, specifically 1:8:28:56:70:56:28:8:1, is influenced by several factors, including the number of unpaired electrons, the spin multiplicity, and the hyperfine coupling with neighboring hydrogen nuclei. Each hydrogen atom in the t-Bu radical contributes to the overall intensity distribution through its possible spin states, leading to a binomial distribution of line intensities that reflects the different combinations of aligned and anti-aligned spins. This specific ratio is significant as it provides insights into the radical's electronic structure and stability, revealing how the distribution of unpaired electrons affects its reactivity and interactions with other molecules. Understanding these ratios enhances our comprehension of radical behavior, particularly in mechanisms of chemical reactions and in the context of radical stabilization.
\\ \textit{Note:} correct because it accurately describes how the number of unpaired electrons and their interactions with neighboring hydrogen nuclei lead to a binomial distribution of line intensities in the EPR spectrum of the t-Bu radical, reflecting the contributions of different spin states to the overall intensity ratios.
\item \textbf{Answer:} Nuclear Magnetic Resonance (NMR) frequency is significant because it directly relates to the nuclear spin properties of isotopes, which influence how they interact with an external magnetic field. The frequency at which a nucleus resonates is determined by its magnetic moment and the strength of the magnetic field applied. For example, 17 O has a nuclear spin of 5/2 and a specific gyromagnetic ratio that results in a resonance frequency of 115.5 MHz when subjected to a 20.0 T magnetic field. This frequency is calculated using the Larmor equation, which links the magnetic field strength to the precession frequency of the nuclear spins, thus highlighting the unique properties of 17 O in NMR studies. Understanding these frequencies is crucial for interpreting NMR spectra and gaining insights into molecular structures and dynamics.
\\ \textit{Note:} correct because it accurately describes how the NMR frequency is influenced by the nuclear spin properties and gyromagnetic ratio of 17 O in relation to the strength of the applied magnetic field, resulting in a specific resonance frequency of 115.5 MHz at 20.0 T.
\end{enumerate}

\vfill
\noindent\textit{End of Answer Key}
\end{document}